\section{E6: the tasks whose ceilings are known}
\label{sec:synthetic}

CIFAR-10 cannot separate ``this repair does not work'' from ``this task does not reward a
second layer''. The first report built two tasks where it can, and their value here is that
the \emph{oracle} row says how much headroom a better training scheme actually has.

\lead{\code{xor}} --- a $32\times32$ Boolean image contains a square and/or a cross; the label
is (square present) XOR (cross present). One layer computes a linear threshold over globally
OR-pooled clause bits, and XOR is not linearly separable in those bits, so it is capped at
$75\%$. Both stacked variants can express it.

\lead{\code{near}} --- both shapes are always present, on a ring of 16 equally spaced slots,
and the label is whether they are close. Rotational symmetry makes each shape's marginal
position identical across classes, so only the \emph{joint} separation carries the label. One
layer and the global-OR stack are at chance by construction; only the spatial stack can
express it.

Layer 1 uses $\pOne\times\pOne$ patches, so no single patch ever contains both shapes.

\newcommand{\syntable}[3]{%
\begin{table}[h]
\centering\small
\caption{#2}
\label{#3}
\pgfplotstabletypeset[
  col sep=comma, string type,
  every head row/.style={before row=\ttoprule, after row=\ttmidrule},
  every last row/.style={after row=\ttbotrule},
  columns={label,fire,dead,l2size,l2pos,acc,sd},
  columns/label/.style={column name=\thead{arm}, column type={l}},
  columns/fire/.style={column name=\thead{L1 fire}, column type={r}},
  columns/dead/.style={column name=\thead{L1 dead}, column type={r}},
  columns/l2size/.style={column name=\thead{L2 size}, column type={r}},
  columns/l2pos/.style={column name=\thead{L2 pos}, column type={r}},
  columns/acc/.style={column name=\thead{test acc (\%)}, column type={r}},
  columns/sd/.style={column name=\thead{$\pm$}, column type={r}},
]{#1}
\end{table}}

\syntable{data/syn_near.csv}{\code{near}: one layer and the global-OR stack are at chance by
construction. \code{L1 fire} is the median per-channel firing rate of the pooled map in
percent and \code{L1 dead} the percentage of channels that never fire; \code{L2 size} and
\code{L2 pos} are the median total and \emph{positive} literal counts of a layer-2 clause ---
what the second layer actually learned.}{tab:synnear}

\syntable{data/syn_xor.csv}{\code{xor}: one layer is capped at $75\%$ by a linear-separability
argument; both stacked variants can express the label. Columns as in
Table~\ref{tab:synnear}.}{tab:synxor}

\begin{figure}[h]
\centering
% Read the table once so the bar labels and the bar heights come from the same file; the
% inline `yticklabels from table={<file>}{...}` form does not inherit `col sep=comma`.
\pgfplotstableread[col sep=comma, header=true]{data/syn_near.csv}\synneartable
\begin{tikzpicture}
\begin{axis}[
  width=0.70\textwidth, height=7.4cm,
  xbar, xmin=45, xmax=104,
  xlabel={\code{near} test accuracy (\%)},
  y dir=reverse, ytick=data, enlarge y limits=0.06,
  yticklabels from table={\synneartable}{label},
  y tick label style={font=\sffamily\scriptsize, align=right},
  bar width=7pt,
  nodes near coords, nodes near coords align={horizontal},
  every node near coord/.append style={font=\sffamily\scriptsize, color=ttgray,
    /pgf/number format/fixed, /pgf/number format/precision=1},
  grid=major, grid style={ttline!40},
  label style={font=\sffamily\small},
  tick label style={font=\sffamily\footnotesize},
  ytick style={draw=none},
]
\addplot[fill=ttorange!75, draw=ttdeep] table[x=acc, y expr=\coordindex] {\synneartable};
\draw[ttslate, dashed, thick] (axis cs:50,-1) -- (axis cs:50,20);
\end{axis}
\end{tikzpicture}
\caption{\code{near}, where a single layer and the global-OR stack are at chance (dashed) by
construction and only the spatial stack can express the label. The oracle row is the first
report's hand-built square/cross detectors.}
\label{fig:near}
\end{figure}

\subsection{\texorpdfstring{\code{near}}{near}: the repairs work, completely}
\label{sec:near}

\begin{ttkey}
On the task that isolates what the second layer is \emph{for}, repairing the first layer's
density solves it outright. The autoencoder layer (\ideaA) reaches
\synNearAutoAuto\,$\pm$\,\synsdNearAutoAuto{}\% and the density controller applied to the
greedy layer (\ideaB) reaches \synNearCalibTask\,$\pm$\,\synsdNearCalibTask{}\%, against
\synNearMctmTask{}\% for greedy training, \synNearMctmOracleSixFour{}\% for the first report's
hand-built oracle, and \synNearSingleTask{}\% --- chance --- for a single layer and for the
global-OR stack, both of which are at chance \emph{by construction}.
\end{ttkey}

Three things follow, and the firing-rate column in Table~\ref{tab:synnear} carries all of them.

\lead{The first report's diagnosis was right, and this is the measurement that shows it.} A
greedy layer~1 on this task is nearly dead --- \syndeadNearMctmTask{}\% of its channels never
fire and the median channel fires on \synfireNearMctmTask{}\% of positions --- because it is
being trained to classify a task one layer cannot classify at all. It still reaches
\synNearMctmTask{}\%, on whatever few channels survive. Move that same layer to
\synfireNearCalibTask{}\% with the controller, changing nothing else and training nothing, and
the stack gains more than ten points.

\lead{The oracle failed for the reason the diagnosis predicts.} The first report built exact
square and cross detectors by hand, and the stack given them did \emph{worse} than the stack
given greedily trained clauses --- which looked like a paradox. It is a density result: an
exact $4\times4$ shape detector fires on a handful of positions in the whole image, and a
layer-2 clause has to find two of them in the same window. Perfect features that never fire
are not usable features. An autoencoder layer is not exact and fires on
\synfireNearAutoAuto{}\% of positions, and it is the one that solves the task.

\lead{Density is necessary and still not sufficient.} Random clauses sit at
\synfireNearRandomRandom{}\% --- comfortably dense --- and reach \synNearRandomRandom{}\%, well
short of the autoencoder's \synNearAutoAuto{}\%. So the ordering on this task is not simply by
firing rate. What the autoencoder adds on top of density is features that are \emph{about} the
shapes, which is exactly what its objective asks for and exactly what a classifier trained on
an unlearnable task cannot discover.

\subsection{\texorpdfstring{\code{xor}}{xor}: the same winner, and a limit on the controller}
\label{sec:xor}

\code{xor} has a different shape: one layer is capped at $75\%$ rather than at chance, and both
stacked variants can express the label, so the question is only whether training finds it.

\ideaA\ finds it. The autoencoder stack reaches \synXorAutoAuto\,$\pm$\,\synsdXorAutoAuto{}\%
--- clear of the one-layer ceiling, and clear of the \synXorMctmOracleSixFour{}\% the
hand-built oracle reaches --- against \synXorMctmTask{}\% for greedy training. Its layer-2
clauses are \synltwoXorAutoAuto\ literals long, of which \synltwoposXorAutoAuto\ are positive:
a composition, not a negation soup, exactly as on \code{near}. Random clauses are second at
\synXorRandomRandom{}\%, also above the oracle --- the same ordering CIFAR-10 gives, and for
the same reason: they are short and they fire.

\ideaB\ does not, and the reason is a real limit on the controller rather than on the idea.
\begin{ttkey}
On \code{xor} the calibrated greedy layer ends up firing on \synfireXorCalibTask{}\% of
positions --- it was asked for $5\%$. These images are almost entirely empty, so a clause of
negated literals matches nearly every patch and a single added pixel literal drops its rate to
almost nothing: there is no single-literal step that lands in the band, and the greedy walk
ends up outside it in one direction or the other. The controller's step size is fine on dense
thermometer-encoded CIFAR-10 patches and too coarse on sparse binary images.
\end{ttkey}

The arm still improves on greedy training (\synXorCalibTask{}\% against \synXorMctmTask{}\%),
and combining it with the autoencoder layer \emph{costs} accuracy
(\synXorAutoCalibAuto{}\% against \synXorAutoAuto{}\%) --- the controller pulls a layer that
was already in a good place out of it. A per-clause step that can adjust more than one literal
at a time, or a band defined in log-odds rather than in rate, would be the obvious fix and is
not tested here.

\subsection{What the second layer becomes}
\label{sec:l2learned}

The proposal's claim about layer~2 was that it would learn a \emph{composition}: a short
conjunction over named channels at named offsets. The first report measured the opposite ---
a thousand-literal clause that is \negMctm{}\% negations and asserts only absences --- and
attributed it to the sparsity of its input. Tables~\ref{tab:synnear} and \ref{tab:synxor}
close that argument, and they do it on both tasks.

\begin{ttkey}
With a repaired layer~1, layer~2 stops being a negation soup and becomes the rule the proposal
described. Over greedy features a layer-2 clause carries \synltwoNearMctmTask\ literals; over
autoencoder features it carries \synltwoNearAutoAuto, of which \synltwoposNearAutoAuto\ are
positive; over calibrated greedy features, \synltwoNearCalibTask\ literals with
\synltwoposNearCalibTask\ positive --- which is to say ``channel $a$ fires here \textsc{and}
channel $b$ fires there'', in as many symbols as it takes to write it down. On \code{xor} the
same thing happens --- \synltwoXorMctmTask\ literals over the greedy layer against
\synltwoXorAutoAuto\ over the autoencoder one. Every arm with a short layer-2 clause is an arm
that works, and on \code{near} every arm whose layer-2 clause runs to four figures sits below
\synNearMctmTask{}\%.
\end{ttkey}

A \synltwoNearCalibTask-literal clause is also an interpretable one, which is the property
Tsetlin machines are supposed to be chosen for and the property the first report found the
stack had lost. It is recovered here for free, as a side effect of fixing the input's density.

\subsection{Credit propagation fails here too, and in the same way}

\code{near} is where \S\ref{sec:credit} predicted the credit path might have bandwidth: the
layer-2 clauses are small and they match. It does not help.

From a random layer~1 both credit arms end at chance (\synNearCreditIaRandom{}\% and
\synNearCreditRandom{}\%) with \syndeadNearCreditIaRandom{}\% of layer-1 channels dead ---
the same saturation-then-extinction as on CIFAR-10, at a tenth the scale. Warm-started from
the greedy layer~1 it reaches \synNearCreditWarmTask{}\% against \synNearMctmTask{}\% for the
frozen stack it started from, and damping the step size to $0.1$ changes that to
\synNearCreditWarmPOneTask{}\%: the credit path is not moving the result in either direction.
The one credit arm that does better, \synNearCreditCalTask{}\%, is the one with \ideaB's
controller running inside it --- and it is still below the \synNearCalibTask{}\% that same
controller reaches with layer~1 simply frozen.

Across two very different tasks, then, the credit path either does nothing or does harm, and
what looks like it helping is the controller underneath it. That is worth knowing: the failure
is a property of the feedback rule, not of the setting it was tried in.

